% Transcription — fonds Alexandre Grothendieck, shelfmark 101, batch 1 (pages 1-20).
% Established from the facsimile archives/batches/101/batch-01.pdf (a mirror of
% https://grothendieck.umontpellier.fr/101.pdf), per the skill
% transcribe-grothendieck. The batch file carries no cover sheet: PDF page N
% is page N of the folder; the archivists' pencil numbers bottom-left agree
% wherever legible (« 10 », « 11 », « 12 », « 14 », « 17 », « 18 », « 19 »,
% « 20 »).
%
% Pass: Opus 5.5 (claude-opus-5-5), 2026-09-24 — first pass, unchecked against the pages by a human.
%
% The folder sits in the inventory group "69-89" « Géométrie et topologie
% combinatoire », whose subject does not match this folder's (as for folder
% 100); the folder has its own date, which \dating{} carries verbatim.
% Nothing is concluded here from the group.
%
% Fifteen of the twenty pages are transcribed. Absent:
%   - pages 1, 5 and 16, sheets of buff card blank but for his words top
%     right — covers; their words head the three \section{} (hand.md §5):
%     « Relations entre cohomologie locale et projective (via les cônes) »,
%     « Dualité locale pour Modules non quasi cohérents », « Notes antiques
%     1958 (Caractérisation des Coh. Mac., des Gorenstein etc) »;
%   - pages 9 and 13, typed pages of some other French text (page « - 85 - »,
%     on « faisceaux fins », and page « II - 18 - », on modules of fractions
%     A[S^-1]), scanned upside down, with no mark of his hand. Pages 8 and 12
%     are versos of typescripts too, but are covered with his own notes and
%     are transcribed; the typed text under them is not.
% No page carries other people's text that he annotated, so nothing is
% summarised under the folder-100/71 rule.
%
% What the batch is. Three runs, each opened by its own cover.
%   2-4    « Relations entre cohomologie locale et globale »: the punctured
%          cone X' = Spec(S) - Y over X = Proj(S), and the comparison of
%          Ext with supports on the cone with graded sums of Ext_T(F, G(n))
%          on X (boxed formula p. 2); passage to the affine cone and to
%          completions, two pairs of spectral sequences (p. 4). Rapid hand,
%          much struck, many connecting words lost.
%   6-15   « Dualité locale pour Modules non quasi cohérents »: p. 6, Ext on
%          Spf A computed by Ext_A(M, N) (Prop. 1, Cor. 2, two proofs); p. 7,
%          Ext_Y on a formal prescheme, broken off; p. 8, scattered notes in
%          black ink (duality for group schemes with values in Q/Z, and
%          pairings H^i_Y(M) x Ext(M, A) -> I in dimension 1); pp. 10-15,
%          local duality Ext^{n-i}(X; F, O_X) -> Hom(H^i_x(F), I) for an
%          arbitrary O_X-Module F on Spec of a complete regular local ring:
%          reduction to F = i_!(O_U), the case n = 1 (triples (M_0, M_1, u),
%          facts a)-c), Ext^1_A(K, A) = 0 for A complete), scratch on
%          adèles and idèle classes (p. 12), and the case n = 2 (pp. 14-15),
%          where phi: A_f -> D(H^2_x(F)) is injective but not surjective —
%          one must take topological duals.
%   17-20  « Notes antiques 1958 »: p. 17 a plan (§ 1 the functors H^i_C,
%          § 2 dualising functors and modules, local duality, Gorenstein
%          rings, § 3 the functors H_i and dualising complexes); p. 18 the
%          Definition of a « module fondamental » Omega (finite type,
%          D(Omega) = H^n(A)) and Theorem 4; p. 19 (cancelled by two long
%          diagonals, transcribed) the Complement to Th. 3 and the pairings
%          H^p(M) x Ext^{n-p}(M, Omega) -> H^n(Omega) -> I; p. 20
%          Corollaries 1-3: the pairing is perfect for all p iff A is
%          Cohen-Macaulay of dimension n; Ext^i(k, Omega) = 0 for i != n,
%          Ext^n(k, Omega) = k; Omega Cohen-Macaulay, Ext^i(Omega, Omega) = 0
%          for i != 0, Hom(Omega, Omega) = A.
% Notation fixed for the batch: D( ) is his duality functor, I the dualising
% (injective hull) module H^n_x(O_X); \underline{Ext} his sheaf Ext; H^i(M)
% for the local cohomology H^i_x(M) where he drops the x. Throughout, the
% connective prose is fast and often not recoverable; the formulas are.
\documentclass[11pt,a4paper]{article}
% Le préambule commun à toutes les transcriptions du fonds Grothendieck.
%
% Il tient en une idée : **l'incertitude se compose, elle ne se résout pas.**
% Une transcription qui lisse un mot douteux a détruit la seule chose pour
% laquelle on la faisait — un lecteur doit pouvoir savoir, à chaque endroit,
% ce qui était sur le feuillet et ce qui a été ajouté. Les six macros qui
% suivent sont donc l'appareil critique, et non de la décoration.
%
% Les mêmes commandes sont reconnues par `scripts/render.mjs`, qui produit la
% vue de lecture du site. Une macro ajoutée ici doit l'être aussi là-bas,
% faute de quoi le rendu échouera bruyamment — ce qui est voulu : un rendu
% silencieusement faux est pire qu'un rendu absent.

\ProvidesPackage{grothendieck}[2026/08/08 Transcriptions du fonds Grothendieck]

\usepackage{amsmath,amssymb,amsthm}
\usepackage{mathrsfs}
\usepackage{stmaryrd}
% Les diagrammes commutatifs sont partout dans ce fonds : c'est la langue
% même de la géométrie algébrique de Grothendieck, et une transcription qui
% les rendrait en prose ne transcrirait rien.
\usepackage{tikz-cd}
% Français par défaut — c'est la langue des feuillets — mais l'anglais est
% chargé aussi : la traduction et le résumé sont des documents anglais, et la
% typographie française (espaces avant les deux-points) y serait une faute.
% Ces documents basculent par \selectlanguage{english} en tête de corps.
\usepackage[english,french]{babel}
\usepackage{fontspec}
\usepackage{geometry}
\geometry{a4paper,margin=25mm,marginparwidth=28mm}
\usepackage{marginnote}
\usepackage{xcolor}
% ulem, not soul. `soul`'s \st and \ul are built on a character-by-character
% scan that XeLaTeX's font machinery breaks: the document fails with an
% undefined control sequence rather than degrading. ulem does the same two jobs
% and is Unicode-engine safe. `normalem` stops it hijacking \emph, which the
% transcriptions use for Grothendieck's own emphasis.
\usepackage[normalem]{ulem}
\usepackage{hyperref}
\hypersetup{colorlinks=true,linkcolor=black,urlcolor=black,pdfborder={0 0 0}}

% --- Métadonnées du lot -----------------------------------------------------
% Reprises telles quelles par le script de rendu, qui les affiche en tête de
% la vue de lecture. Elles fixent ce que le fichier prétend couvrir : sans
% elles, une transcription flotte sans point d'attache dans un fonds de seize
% mille feuillets.

\newcommand{\folder}[1]{\gdef\grothFolder{#1}}
\newcommand{\batch}[1]{\gdef\grothBatch{#1}}
\newcommand{\leaves}[2]{\gdef\grothFirst{#1}\gdef\grothLast{#2}}
\newcommand{\foldertitle}[1]{\gdef\grothFoldertitle{#1}}
\gdef\grothFolder{?}\gdef\grothBatch{?}\gdef\grothFirst{?}\gdef\grothLast{?}\gdef\grothFoldertitle{}

% --- Le résumé --------------------------------------------------------------
%
% Il ouvre la lecture modernisée, sous le titre « Résumé », et il est composé
% en petit. Ce n'est pas un ornement : c'est le seul passage du document qui
% ne soit pas des mathématiques, et un lecteur doit pouvoir le reconnaître
% avant de l'avoir lu — puis le sauter s'il connaît déjà le sujet, ou s'y
% arrêter s'il ne le connaît pas. Le filet qui le suit marque la reprise du
% corps.

\newenvironment{resume}
  {\small\setlength{\parskip}{0.45em}}
  {\par\medskip\hrule\medskip}

% --- Mots-clés ---------------------------------------------------------------
%
% En anglais, en clôture du résumé de la lecture modernisée : le vocabulaire
% moderne sous lequel chercher ce que les feuillets construisent. Le manifeste
% du site (`npm run manifest`) les extrait pour étiqueter la cote entière —
% cette ligne est la source unique des tags, il n'y a pas de fichier à côté.
\newcommand{\keywords}[1]{\par\smallskip\noindent
  {\footnotesize\textsc{Keywords} --- \textit{#1}}\par}

% --- L'appareil critique ----------------------------------------------------

% Le numéro de feuillet, en marge. C'est l'ancre qui permet de régler un
% désaccord sur une formule en regardant *un* feuillet plutôt qu'une chemise
% de six cents. Le site s'en sert aussi pour tourner les pages du fac-similé
% quand on fait défiler la transcription.
\newcommand{\leaf}[1]{%
  \marginnote{\footnotesize\color{gray!70}\textsf{#1}}%
  \label{leaf:#1}\ignorespaces}

% Illisible : on ne devine pas. Un mot inventé qui se lit comme les autres est
% le pire résultat possible.
\newcommand{\ill}{\textcolor{red!60!black}{[\,\ldots\,]}}

% Lecture douteuse : le mot est donné, et signalé comme incertain.
\newcommand{\uncertain}[1]{\uline{#1}}

% Ajout de l'éditeur — une lettre manquante, un mot sous-entendu.
\newcommand{\add}[1]{\textcolor{blue!50!black}{[#1]}}

% Biffé par Grothendieck lui-même. Ce qu'il a rayé fait partie du document :
% une rature dit souvent où la pensée a changé de direction.
\newcommand{\struck}[1]{\sout{#1}}

% Note du transcripteur, sur l'état matériel du feuillet.
\newcommand{\note}[1]{\par\smallskip{\small\color{gray!60!black}\textit{[#1]}}\par\smallskip}

% Note marginale de Grothendieck, distincte du corps du texte.
\newcommand{\marginal}[1]{\par\smallskip\noindent\hspace{1em}%
  {\small\color{gray!60!black}#1}\par\smallskip}

% --- Titre ------------------------------------------------------------------

\AtBeginDocument{%
  \begin{center}
    {\large\bfseries Fonds Alexandre Grothendieck --- cote n\textsuperscript{o}~\grothFolder}\\[2pt]
    {\itshape\grothFoldertitle}\\[4pt]
    {\small feuillets \grothFirst--\grothLast\ (lot \grothBatch)}\\[2pt]
    {\footnotesize\color{gray!60!black}Transcription établie d'après le fac-similé numérisé
     par l'Université de Montpellier.\\
     Les lectures incertaines sont soulignées, les illisibles notées [\,\ldots\,],
     les ajouts entre crochets.}
  \end{center}
  \vspace{1em}}

\endinput


\folder{101}
\batch{1}
\pages{1}{20}
\dating{[à partir de 1967-1968]}
\watermark{Édition de démonstration}
\foldertitle{Cohomologie locale : notes manuscrites (s.d.).}

\begin{document}

\section{Relations entre cohomologie locale et projective (via les cônes)}
\note{titre de sa main, en haut à droite du feuillet 1, une chemise de papier
bistre par ailleurs vierge, qui ne reçoit pas de numéro de page}

\page{2}
\textbf{Relations entre cohomologie locale et globale}
\note{titre souligné, en tête du feuillet ; il diffère de celui de la chemise
(« projective » y est remplacé par « globale »)}

$Y$ préschéma de base, $\mathcal{S}$ Algèbre graduée quasi-cohérente sur $Y$,
à degrés $\geqslant 0$, engendrée par $\mathcal{S}_1$ de type fini,
$X = \operatorname{Proj}(\mathcal{S})$, $X' = \operatorname{Spec}(\mathcal{S})
- Y$ le cône épointé, $\overline{X'}$ le cône affine. Donc
\[
  X' \simeq \mathbf{V}(\mathcal{O}_X(1)) - \sigma(X)
  = \operatorname{Spec} \struck{\ill{}} \Bigl(\coprod_{n \in \mathbf{Z}}
  \mathcal{O}(-n)\Bigr).
\]
\note{le début de la ligne, avant « $\mathbf{V}$ », est d'une lecture
incertaine ; le mot biffé après « $\operatorname{Spec}$ » commence peut-être
par « Sym »}

Soit $Z$ une partie fermée de $X$, $Z'$ son image inverse dans $X'$,
$\overline{Z'}$ son adhérence [définie aussi au moyen de l'Idéal homogène
\uncertain{définissant} $Z$]. $U$ \uncertain{un ouvert} de $X$, $U'$ son image
inverse dans $X'$, $\overline{U'}$ dans $\overline{X'}$. Soient $F$, $G$ deux
Modules [quasi-cohérents] sur \struck{$U$} $X$, $F'$, $G'$ \ill{} \uncertain{les
images inverses} sur $U'$.
\note{le début de la plupart des lignes de ce paragraphe, au bord gauche, est
d'une lecture incertaine ; les crochets sont de lui}

\[
\begin{aligned}
  \operatorname{Ext}^{\bullet}_{Z' \cap U' = T'}(X'; F'; G')
  &= H^{\bullet}\, \struck{\ill{}}\; \underline{\operatorname{Hom}}_{Z' \cap U'\,
  T'}(F'; C(G'))\\
  &= H^{\bullet}\, \struck{\ill{}}\; \underline{\operatorname{Hom}}_{Z \cap U\,
  T}(F, p_{U*} C(G'))\\
  &= H^{\bullet}\, \underline{\operatorname{Hom}}_{Z \cap U\, T}(F,
  C(p_{U*} G'))\\
  &= \operatorname{Ext}^{\bullet}_{Z \cap U\, T}\Bigl(F, \coprod_{n \in
  \mathbf{Z}} G(n)\Bigr)
\end{aligned}
\]
\marginal{$G$ q.-coh.}
\marginal{$G$ q.-coh., $F$ de type fini}
\note{les deux marginales sont reliées par des accolades aux deux dernières
égalités ; les indices « $Z' \cap U'$ » et « $Z \cap U$ » sont soulignés par
une accolade horizontale qui les renomme $T'$ et $T$ ; deux passages de la
première et de la deuxième ligne sont biffés en hachures serrées et ne se
lisent pas}

À gauche, en regard :
\[
  T = U \cap Z \ \text{(loc. fermé dans $X$)}, \qquad T' = U' \cap Z',
\]
\[
  p : X' \longrightarrow X, \qquad p_U : U' \longrightarrow U.
\]

Encadré :
\[
  \operatorname{Ext}^i_{\struck{\ill{}}\,T',\, \mathcal{O}_{X'}}(X'; F', G')
  \simeq \coprod_{n \in \mathbf{Z}} \operatorname{Ext}^i_{T\,
  \struck{\ill{}}}(F, G(n)).
\]
\note{l'indice de gauche est surchargé ; « $\mathcal{O}_{X'}$ » est écrit
dessous d'une encre plus claire}

\page{3}
Localisons sur $Y$, \ill{} \ill{}.
\[
  \underline{\operatorname{Ext}}^{\bullet}_{\overline{T'}\,
  \struck{\ill{}},\, \mathcal{O}_{\overline{X'}}}(\overline{X'}/Y;
  \overline{F'}, \overline{G'}) \simeq \coprod_{n \in \mathbf{Z}}
  \underline{\operatorname{Ext}}^{i}_{\struck{\ill{}},\, \mathcal{O}_U}(X/Y;
  F, G(n))
\]
\note{les indices sont en partie biffés ; au second membre on lit
« $\struck{X}\,X/Y$ » ; une accolade isolée, en dessous, porte
« $\underline{\operatorname{Ext}}_{Z' \cap U'}$ »}

\marginal{N.B. \ill{} \ill{} \ill{}\,…}
\note{note cerclée dans la marge gauche, écrite en biais, que nous ne lisons
pas au-delà de « N.B. » ; au-dessous, un croquis du cône : la droite
$\overline{Z'}$, le cône $\overline{X'}$ et, cerclé, le sommet au-dessus
de $y$}

\struck{D'autre part, on aura \ill{} $W$ \ill{} pour simplifier $W = X$, donc
$U' = X'$, $\overline{U'} = \overline{X'}$.}
\struck{$\operatorname{Ext}^{\bullet}_{Z', \mathcal{O}_{X'}}(X'; F', G')
\Leftarrow H^p(X', \underline{\operatorname{Ext}}^q_{Z'}(F', G'))$}
\note{ces deux lignes sont barrées de traits obliques ; la dernière porte une
flèche de suite spectrale}

\[
  \operatorname{Ext}^{\bullet}_{\overline{T'}}(\overline{X'}; \overline{F'},
  \overline{G'})
  \ \Big\uparrow\
  \operatorname{Ext}^{\bullet}_{\widehat{T'}}(\widehat{X'}; \widehat{F'},
  \widehat{G'})
\]
\note{la flèche verticale va du second vers le premier ; à gauche, deux
croquis : la droite $\overline{Z'}$ coupant $\overline{X'} - \overline{U'}$
en un point marqué, puis $\overline{Z'}$ privé d'un point}

\[
  H^p\bigl(\overline{Z'} - \overline{R'},\ \underline{\operatorname{Ext}}^q_{\overline{Z'}}
  (\overline{F'}, \overline{G'})\bigr)
\]
\marginal{\uncertain{Justifier} \ill{} \ill{} \ill{} \uncertain{compatibilité}
\ill{} \uncertain{extension} \ill{} \ill{}\,…}
\note{note cerclée, rattachée par un crochet à « $\overline{Z'} -
\overline{R'}$ »}

\[
  H^p(\overline{Z'} - \overline{R'}, \overline{E})
  \overset{(2)}{\longleftrightarrow}
  H^p(\widehat{Z'} - \widehat{R'}, \widehat{E})
\]
\marginal{bij.-bij. \uncertain{par localisation}}
\note{à gauche, deux croquis de plans : l'un traversé par une droite épaisse
portant un point, l'autre marqué d'un seul point}

Le \uncertain{lemme} des \uncertain{cinq} \ill{} \ill{}
$H^p_{\overline{Z'}}(\overline{R'}, \overline{E})$, \ill{} \ill{}
\uncertain{donnera} \ill{} \ill{} $H^p_{\overline{Z'}}(\overline{X'},
\overline{E})$.

\struck{Mais ceci} $\longrightarrow$ [Dégageons un \emph{lemme} \ill{} les
données de base \uncertain{plutôt} \ill{} \ill{} \ill{}
$\underline{\operatorname{Ext}}^{\bullet}_T(X, F, G)$.]
\note{ce crochet est écrit à droite, au bas du feuillet, et ne se lit qu'en
partie}

\page{4}
\[
\begin{aligned}
  \operatorname{Ext}^{\bullet}_{T, \mathcal{O}_X}(X, F, G(n))
  &\Longleftarrow H^p\bigl(U, \underline{\operatorname{Ext}}^q_{T,
  \mathcal{O}_U}(F, G(n))\bigr)\\
  \operatorname{Ext}_{T', \mathcal{O}_{X'}}(\overline{X'}; \overline{F'},
  \overline{G'})
  &\Longleftarrow H^p\bigl(U', \underline{\operatorname{Ext}}_{T',
  \mathcal{O}_{U'}}(\overline{F'}, \overline{G'})\bigr)
\end{aligned}
\]
\note{les deux lignes sont reliées, terme à terme, par des signes « $=$ »
verticaux}

[\uncertain{Compatibles} \ill{} \uncertain{formations} des
$\underline{\operatorname{Ext}}^q_{T}(F, G)$ et
$\underline{\operatorname{Ext}}^q_{T'}(F', G')$]
\note{l'indice du premier $\underline{\operatorname{Ext}}$ est surchargé
d'une tache}

\[
\begin{aligned}
  \operatorname{Ext}^{\bullet}_{T, \mathcal{O}_X}(X; F, G(n))
  &\Longleftarrow H^p_T\bigl(\struck{X}\, X,
  \underline{\operatorname{Ext}}^q(F, G)(n)\bigr)\\
  \operatorname{Ext}_{T', \mathcal{O}_{X'}}(\overline{X'}, \overline{F'},
  \overline{G'})
  &\Longleftarrow H^p_{T'}\bigl(\overline{X'},
  \underline{\operatorname{Ext}}^q(F', G')(n)\bigr)
\end{aligned}
\]
\note{sous le premier membre de la première ligne, un signe « $\vee$ » ;
sous le second, un « $=$ » vertical}

\marginal{N.B. les deux \uncertain{suites spectrales} \ill{} \ill{} \ill{}
\uncertain{restrictions} \ill{} \ill{} \ill{} \uncertain{graduées}}
\note{note de la marge gauche, écrite en biais, reliée par deux accolades aux
deux paires de suites spectrales}

\uncertain{Cas particuliers} lorsque $F = \mathcal{O}_X$\,…
\note{le feuillet s'arrête sur ces points de suspension}

\section{Dualité locale pour Modules non quasi cohérents}
\note{titre de sa main, en haut à droite du feuillet 5, une chemise de papier
bistre par ailleurs vierge, qui ne reçoit pas de numéro de page ; « non quasi
cohérents » est souligné}

\page{6}
\textbf{Proposition 1.} Soit $A$ un anneau \uncertain{adique} noethérien.
\ill{} \uncertain{supposé} \uncertain{séparé} \uncertain{complet}, $M$,
\struck{$N$} $N$ deux $A$-modules de type fini, et $\mathfrak{X} =
\operatorname{Spf} A$, $\mathcal{M}$, $\mathcal{N}$ les Modules cohérents sur
$\mathfrak{X}$ définis par $M$, $N$. Alors on a des isomorphismes canoniques
\[
  \operatorname{Ext}^i(\mathfrak{X}; \mathcal{M}, \mathcal{N}) \simeq
  \operatorname{Ext}^i_A(M, N).
\]
\note{écriture rapide ; la fin de la première ligne est d'une lecture très
incertaine}

\textbf{Corollaire 2.} On a
\[
  \underline{\operatorname{Ext}}^i(\mathcal{M}, \mathcal{N}) \simeq
  \operatorname{Ext}^i_A(M, N)^{\Delta}\,\struck{\ill{}}.
\]

[En effet, on sait que les $\underline{\operatorname{Ext}}^i(\mathcal{M},
\mathcal{N})$ sont cohérents, donc \ill{} \uncertain{par} la nullité \ill{}
\ill{} de la suite spectrale \ill{} \ill{} cohérents sur $\mathfrak{X}$,
\ill{} \ill{}
\[
  \operatorname{Ext}^i_{\mathfrak{X}}(\mathcal{M}, \mathcal{N}) \simeq
  H^0\bigl(\mathfrak{X}, \underline{\operatorname{Ext}}^i(\mathcal{M},
  \mathcal{N})\bigr)
\]
\struck{et comme $\operatorname{Ext}^i(\mathfrak{X}, \mathcal{M}, \mathcal{N})
\simeq \operatorname{Ext}^i_A(M, N)$, \ill{} \ill{}}
d'où
\[
  H^0\bigl(\mathfrak{X}, \underline{\operatorname{Ext}}^i(\mathcal{M},
  \mathcal{N})\bigr) \simeq \operatorname{Ext}^i_A(M, N) \quad \text{pour
  tout } i,
\]
d'où le cor.~2 comme $\underline{\operatorname{Ext}}^i(\mathcal{M},
\mathcal{N})$ est \uncertain{cohérent}.]
\note{ce crochet est une démonstration de la proposition 1 par le corollaire,
qu'il reprend ensuite dans l'autre sens}

Il suffit en effet de prendre une résolution \struck{finie} de $M$ par des
modules libres de type fini \uncertain{pour} \uncertain{démontrer} le
résultat. \uncertain{Autre} procédé : comme $\mathfrak{X} \to X =
\operatorname{Spec} A$ est plat, les foncteurs \ill{} \uncertain{commutent} à
la formation des $\underline{\operatorname{Ext}}^i$, d'où
\[
  \underline{\operatorname{Ext}}^i_{\mathcal{O}_{\mathfrak{X}}}(\mathcal{M},
  \mathcal{N}) \simeq
  \underline{\operatorname{Ext}}^i_{\mathcal{O}_X}(\widetilde{M},
  \widetilde{N})^{\wedge} \quad \text{\ill{}}\ \
  \underline{\operatorname{Ext}}^i(\widetilde{M}, \widetilde{N})
  = \operatorname{Ext}^i_A(M, N)^{\sim} \ \ (\text{\uncertain{cq}}),
\]
d'où
\[
  \underline{\operatorname{Ext}}^i_{\mathcal{O}_{\mathfrak{X}}}(\mathcal{M},
  \mathcal{N}) \simeq \operatorname{Ext}^i_A(M, N)^{\Delta}
\]
(c'est le cor.~2) et en prenant les sections, on trouve à nouveau la
prop.~1.

\page{7}
Soit \struck{\ill{}} $\mathfrak{X}$ un préschéma formel localement noethérien,
$Y$ une partie fermée de $\mathfrak{X}$, $\mathcal{F}$, $\mathcal{G}$ deux
Modules \ill{} \uncertain{cohérents} sur $\mathfrak{X}$, \ill{} \ill{}
\ill{} \ill{}
\[
  \operatorname{Ext}^i_Y(\mathfrak{X}; \mathcal{F}, \mathcal{G})
  \Longleftarrow H^p_Y\bigl(\mathfrak{X},
  \underline{\operatorname{Ext}}^q(\mathcal{F}, \mathcal{G})\bigr),
\]
en particulier \ill{} \ill{} \ill{}, \ill{} \ill{} \ill{} lesquels il
\uncertain{faut} \uncertain{préciser} \ill{} \ill{} \ill{} limite projective
\[
  \varprojlim_n \operatorname{Ext}^i_Y(\mathfrak{X}, \mathcal{F},
  \mathcal{G}_n).
\]
Notons d'ailleurs que $\operatorname{Ext}^{*}_Y(\mathfrak{X}, \mathcal{F},
\mathcal{G}_n)$
\note{la phrase s'arrête là ; le reste du feuillet est blanc}

\page{8}
\note{notes éparses à l'encre noire, sur le verso d'un tapuscrit
sans rapport}
\struck{$E(\mathcal{U},$} \quad \struck{$H^i$}
\[
\left\{
\begin{aligned}
  \underline{E}_{U/k}(G, {}_{\infty}\mu) &\simeq \underline{E}_k(N^i_{U/k}\,G,
  \mathbf{Q}/\mathbf{Z})\\
  \underline{E}_{S/k}(G, {}_{\infty}\mu) &\simeq
  \underline{E}_k(N^i_{Y/k}(G), \mathbf{Q}/\mathbf{Z})\\
  \underline{E}_{Y/k}(G, {}_{\infty}\mu) &\simeq E_k(N^i_{S/k}(G),
  \mathbf{Q}/\mathbf{Z}).
\end{aligned}
\right.
\]
\note{les indices $Y/k$ et $S/k$ des deux dernières lignes se croisent
ainsi sur la page ; nous les laissons tels quels}

\struck{$\operatorname{Ext}^{*}_K(G, K)$}

$Y \subset S$
\[
\begin{array}{ccc}
  0 & & \\
  \downarrow & & \\
  H^0_Y(M) & \text{---} & \operatorname{Ext}^1_Y(S; M, A)\\
  \downarrow & & \uparrow\\
  H^0(S, M) & \text{---} & \operatorname{Ext}^i_Y(S; M, A)\\
  \downarrow & & \\
  H^0(U, M) & &
\end{array}
\]
\note{deux colonnes accouplées terme à terme par des tirets ; l'exposant
« $1$ » de la première ligne de droite est surchargé}

$k \leftarrow V \rightarrow K$

\[
\begin{array}{ccccc}
  0 & & & & \\
  \downarrow & & & & \\
  H^0(M) & \times & \operatorname{Ext}^{\ill{}}(M, A) & \longrightarrow &
  I\\
  \downarrow & & \uparrow & & \\
  M & \times & \operatorname{Ext}^{\ill{}}(M, A) & \longrightarrow & I\\
  \downarrow & & \uparrow & & \\
  M \otimes K & \times & \operatorname{Hom}(M_K, K) & \longrightarrow & I\\
  \downarrow & & \uparrow & & \\
  H^1(M) & \times & \operatorname{Ext}^0(M, A) & \longrightarrow & I\\
  \downarrow & & \uparrow & & \\
  0 & & 0 & &
\end{array}
\]
\note{à droite de la troisième ligne : « \uncertain{Hom}$(K, I) \simeq K$ »
; sous le tableau, isolé, « $\operatorname{Ext}^i$ » ; les exposants des
deux premiers $\operatorname{Ext}$ ne se lisent pas, et le « $M \otimes K$ »
est surchargé}

\[
  \operatorname{Ext}_Y(M, A) \simeq \operatorname{Hom}(H^i_Y(X, M), I)
\]

\page{10}
\textbf{\emph{Dualité locale.}} $A$ anneau local noethérien \emph{complet}
\emph{régulier} de dimension $n$, $X = \operatorname{Spec} A$, $x$ son point
fermé, $F$ Module variable sur $X$ (\emph{pas nécessairement
quasi-cohérent}\,!)
\[
  H^i_x(F) \times \operatorname{Ext}^{n-i}(X; F, \mathcal{O}_X)
  \longrightarrow H^n_x(\mathcal{O}_X) = I \qquad (I \text{ module
  \uncertain{dualisant}})
\]
d'où
\[
  \boxed{\operatorname{Ext}^{n-i}(X; F, \mathcal{O}_X) \longrightarrow
  \operatorname{Hom}(H^i_x(F), I)}
\]
Je dis que ce sont des \emph{isomorphismes} pour \emph{tout $i$}, \emph{à
\uncertain{priori}} \ill{} \emph{\uncertain{chaque} pris} \ill{}.
\note{« pas nécessairement quasi-cohérent » est doublement souligné ; la fin
de la phrase « Je dis que » est écrite sur deux lignes et soulignée par
morceaux}

On est ramené : la question \ill{} \ill{} pour un faisceau $F$ de la forme
$\mathcal{O}_{U,X} = i_!(\mathcal{O}_U)$, où $U$ est un ouvert \ill{} \ill{}
de $X$ (\uncertain{voire} \ill{} de la forme $X_f$) et $i : U \to X$.
L'isomorphisme \ill{} \uncertain{canonique} \ill{} \ill{} alors
\[
  \operatorname{Ext}^{n-i}(U; \mathcal{O}_X, \mathcal{O}_X)
  \xrightarrow{\ \sim\ } \operatorname{Hom}(H^i_x(\mathcal{O}_{U,X}), I),
\]
i.e.
\[
\left\{
\begin{aligned}
  &H^i_x(\mathcal{O}_{U,X}) = 0 \quad \text{si } i \neq n,\\
  &D\bigl(H^n_x(\mathcal{O}_{U,X})\bigr) \simeq \struck{\operatorname{Hom}}\,
  \Gamma(U, \mathcal{O}_X) = A_f.
\end{aligned}
\right.
\]
\struck{\ill{}} On a \ill{} si $U = X$, \ill{} \ill{}, \uncertain{Reste} $U
\neq X$, \ill{}.

\struck{$H^0_x$}

1) $n = 0$ \ill{} obtient une Tautologie.

2) $n = 1$, il suffit de prendre $U =$ \uncertain{ouvert} \uncertain{principal},
et le faisceau \uncertain{devient}, puisque $H^0_x(\mathcal{O}_{U,X}) = 0$,
$H^1_x(\mathcal{O}_{U,X}) = K$,
\[
  \operatorname{Hom}(K, K/A) \simeq K
\]
(autodualité bien \uncertain{connue}\,!). Notons que \struck{pour} la
catégorie des $\mathcal{O}_X$-Modules est \uncertain{équivalente} à la
catégorie des triplets $(M_0, M_1, u)$, où $M_0$ est un $A$-module, $M_1$ un
$K$-module [resp.\ $F_x$ et $F_{x'}$, où $x'$ est le point générique] et
\[
  u : M_0 \longrightarrow M_1
\]
un $(A \to K)$-homomorphisme. On aura
\[
  H^0_x(F) = \operatorname{Ker} u, \quad H^1_x(F) = \operatorname{Coker} u
  \qquad [H^0(X, F) \simeq M_0,\ H^i(X, F) = 0 \ i \neq 0]
\]
\note{sous le crochet, une ligne biffée qui ne se lit pas}

\page{11}
D'autre part, on a une résolution injective \uncertain{globale} de
$\mathcal{O}_X$ \uncertain{par}
\[
  \mathcal{O}_X \longrightarrow \underbrace{\underline{K}_X \longrightarrow
  \underline{K}_X/\mathcal{O}_X \longrightarrow 0 \longrightarrow 0
  \longrightarrow \cdots}_{K^{\bullet}}
\]
d'où
\[
  \operatorname{Ext}^{*}(X; F, \mathcal{O}_X) \simeq H^{*}\bigl(0
  \longrightarrow \underset{0}{\operatorname{Hom}(M_1, K)} \longrightarrow
  \underset{1}{\operatorname{Hom}(M_0, K/A)} \longrightarrow 0 \cdots\bigr)
\]
i.e.\ nuls en \uncertain{dimension} $\neq 0, 1$ et
\[
\begin{aligned}
  \operatorname{Ext}^0(X; F, \mathcal{O}_X) &\simeq
  \operatorname{Ker}\bigl(\operatorname{Hom}(M_1, K) \to
  \operatorname{Hom}(M_0, K/A)\bigr)\\
  \operatorname{Ext}^1(X; F, \mathcal{O}_X) &\simeq
  \operatorname{Coker}(\quad \text{id} \quad)
\end{aligned}
\]
Dans le \uncertain{cas} \ill{} les \uncertain{homomorphismes} dualité
deviennent
\[
\begin{aligned}
  E^0 &\overset{\varphi^0}{\longrightarrow}
  \operatorname{Hom}(\operatorname{Coker} u, K/A)\\
  E^1 &\overset{\varphi^1}{\longrightarrow}
  \operatorname{Hom}(\operatorname{Ker} u, K/A)
\end{aligned}
\]
\uncertain{Faisons} \uncertain{usage} des faits suivants :

a) Tout homomorphisme $M_1 \to A$ est nul (l'image étant un sous-module
\uncertain{qui est} $K$-\uncertain{module} de $A$) \uncertain{donc}
$\varphi^0$ injectif.

b) Tout homomorphisme $M_1 \to K/A$ se relève en $M_1 \to K$, i.e.\
$\operatorname{Ext}^1_A(M_1, A) = 0$, \uncertain{ou} \uncertain{aussi}
$\underline{\operatorname{Ext}^1_A(K, A) = 0}$ [\uncertain{Pour} le
\uncertain{voir}, \uncertain{utilisons} la résolution injective $A \to$ \ill{}
$A$ \ill{} $\to K \to K/A \to 0$, \ill{} \ill{} \uncertain{exacte}
\[
  0 \longrightarrow \underset{\substack{\wr\\ K}}{\operatorname{Hom}_A(K,
  K)} \longrightarrow \operatorname{Hom}(K, K/A) \longrightarrow 0,
\]
\uncertain{qui} \ill{} \uncertain{utilise} \uncertain{l'autodualité} de
$K$ \ill{}]. \uncertain{Donc}
\note{au-dessus de « $K$ », en fin de crochet, une insertion
« $\widehat{K} = K$ » ; sous « autodualité », une remarque entre parenthèses,
« \uncertain{on} \ill{} \uncertain{fait} \uncertain{que} $A$
\uncertain{complet} », rattachée par un trait}

$\varphi^0$ surjectif.

c) $K/A$ \uncertain{est} injectif, \uncertain{donc} tout homomorphisme
$\operatorname{Ker} u \to K/A$ \uncertain{s'étend} : $M_0 \to K/A$, donc
$\varphi^1$ est surjectif, et tout homomorphisme de $M_0/\operatorname{Ker} u
\subset M_1$ dans $K/A$ \uncertain{s'étend} \ill{} \ill{} $M_1 \to K/A$,
\uncertain{qui} \ill{} \uncertain{vient} \uncertain{donc} \uncertain{d'un}
\ill{} $M_1 \to K$, donc $\varphi^1$ est injectif.
\note{la fin de c) ne se lit qu'en partie ; l'ordre des deux conclusions
« surjectif » et « injectif » pour $\varphi^1$ est celui du feuillet}

\page{12}
\note{feuillet de brouillon, notes éparses à l'encre sur le verso d'un
tapuscrit sans rapport ; nous les donnons dans l'ordre où elles se
présentent, de haut en bas}

$\Phi_x$ \quad $\mathcal{O}_y$ \quad $\mathcal{O}_{X,u}$ \quad $K$
\[
  \operatorname{Ext}^i(K, A) = 0 \qquad \operatorname{Ext}^i(K, A)
\]
$H^0_x$ \quad $H^1_x$
\note{ces premières notes sont isolées par un trait horizontal et un filet
vertical qui les encadrent}

\uncertain{résolution} de $A$
\[
  0 \longrightarrow K \longrightarrow \coprod_{p} K/A_p \longrightarrow I
  \longrightarrow 0
\]
\[
  0 \longrightarrow K \longrightarrow \prod_{p}{}^{\wedge}\, \widehat{K}_p
  \longrightarrow \operatorname{Hom}(K, I) \longrightarrow 0
\]
\note{une accolade au-dessus de la première suite la désigne comme la
résolution de $A$ ; sous « $\prod^{\wedge} \widehat{K}_p$ », entre
guillemets, « produit restreint », puis « i.e.\ \uncertain{adèles} », doublement
souligné ; sous « $\operatorname{Hom}(K, I)$ », un « $=$ » vertical et
« $\varinjlim_{f \in A - 0} f I$ »}

\[
  \operatorname{Ext}^i_A(K, A) \qquad \operatorname{Ext}^i(A,\ \ )
  \qquad \operatorname{Hom}(K, A) = 0
\]
\struck{$\operatorname{Ext}^i$} \quad \struck{$\operatorname{Ext}^i$}
\[
  \operatorname{Ext}^1(K, A) \qquad 0 \longrightarrow A \longrightarrow E
  \longrightarrow K \longrightarrow 0
\]
\[
  0 \longrightarrow A \longrightarrow K \longrightarrow \coprod
  \qquad \operatorname{Hom}(K, A/\mathfrak{m}^n)
\]
\[
  \operatorname{Ext}^1(K,\ \ ) \qquad \operatorname{Ext}^i \qquad
  \operatorname{Ext}^1(\underline{K}, A^{(I)})
\]

\[
\begin{array}{ccc}
  K & \longrightarrow & \prod \widehat{K}_p\\
  \uparrow & & \uparrow\\
  A & \longrightarrow & \prod \widehat{A}_p
\end{array}
\qquad
\left|
\begin{array}{ll}
  0 & \text{Classes d'idèles}\\[4pt]
  1 & \text{Classes d'\uncertain{idéaux}}
\end{array}
\right.
\]
\note{sous chacun des deux produits, un indice abrégé que nous ne lisons
pas ; plus bas, isolé, un « $A$ » souligné}

\page{14}
$n = 2$) \quad $U = X_f$, $f \in \mathfrak{m}_A$, $f \neq 0$, \quad $F =
\mathcal{O}_{X,U}$. Soit $X' = X - x$, \struck{$X$} $Y = V(f)$, $Y' = Y \cap
X'$, $x$ : le point fermé, en \uncertain{conséquence} de $Y$, \ill{}.
\note{« $Y' = Y \cap X'$ » est ajouté en haut à droite, sous un
« $F = \ill{}\,X'$ » biffé}
\[
  H^0_x(F),\ H^1_x(F) = \operatorname{Ker} \text{ et } \operatorname{Coker}
  \text{ de } \underset{\substack{\|\\ 0}}{H^0(X, F)} \longrightarrow
  \underset{\substack{\|\\ 0}}{H^0(X', F)} \quad \text{\uncertain{nuls}}.
\]
$H^i_x(F) = 0$ pour $i > 2$ pour \uncertain{raison} de dimension
\uncertain{cohomologique}.

\uncertain{Reste} : calculer $H^2_x(F)$ et comparer à son dual : $A_f =
H^0(U, \mathcal{O}_U)$.

\uncertain{Comme} \struck{\ill{} \ill{}} pour $H^i(X, F) = 0$ pour \struck{$H$}
$i \neq 0$, $\underline{\forall F}$,
\[
  H^2_x(F) \simeq H^1(X', F).
\]
Soit \struck{$\mathcal{O}_Y$ $F$} $i : Y' \to X'$ l'injection, et $G =
i^{*}\mathcal{O}_X$, \ill{}, \ill{} \uncertain{suite} exacte
\[
  0 \longrightarrow F \longrightarrow \mathcal{O}_{X'} \longrightarrow G
  \longrightarrow 0,
\]
\uncertain{et} \uncertain{cohomologie} \ill{} $Y'$, l'ensemble
\uncertain{fini} \emph{discret} $Y'$. D'où la suite
\uncertain{exacte} \uncertain{longue}
\[
  0 \to \underset{\substack{\|\\ A}}{H^0(X', \mathcal{O}_{X'})}
  \xrightarrow{\alpha^0}
  \underset{\substack{\|\\ \prod \mathcal{O}_{x_i}\\ \|\\ \varinjlim_g A
  g^{-1}}}{H^0(Y'; G)}
  \xrightarrow{\partial}
  \underset{\substack{\|\\ H^2_x(F)}}{H^1(X'; F)}
  \xrightarrow{\alpha^1}
  \underset{\substack{\|\\ I}}{H^1(X'; \mathcal{O}_{X'})} \to 0
\]
\marginal{$g$ \uncertain{variant} \uncertain{dans} \ill{} $x_i$, i.e.\ \ill{}
\ill{} $f$, i.e.\ $(f, g)$ \uncertain{engendrant} \uncertain{un} idéal primaire
pour $\mathfrak{m}$}
\note{la marginale est écrite à gauche, en regard de la limite inductive
sur $g$}

Prenant les \uncertain{duaux} : \ill{} \uncertain{dans} $I$, \ill{}
\uncertain{trouve}

\begin{tikzcd}[column sep=small]
0 & I \arrow[l] & \varprojlim_g gI \arrow[l, "\alpha^{0*}"'] & D(H^2_x(F)) \arrow[l, "\partial^{*}"'] & A \arrow[l, "\alpha^{1*}"'] \arrow[dl, dashed] & 0 \arrow[l] \\
 & & & A_f \arrow[u, "\varphi"] \arrow[ul, dashed, "d"] & &
\end{tikzcd}

\note{sous « $I$ », « $= 1 \cdot I$ » ; sous « $A_f$ », « $=
\varinjlim_m A f^{-m}$ » ; la flèche pointillée de $A$ vers $A_f$ porte
« en $A$-modules »}

L'application $A_f \to \varprojlim_g gI$ s'explicite par
\[
  d_g : h \longmapsto [h]\bigl[\tfrac{1}{g}\bigr].
\]
Le noyau de $d_g$ est donc l'image de $A \to A_f$, donc \ill{}
$\varphi$ \emph{est injection}. \uncertain{Le} \ill{} \ill{} l'image de
$d$ est $\neq \operatorname{Ker} \alpha^{0*}$ ; $\varphi$ \emph{n'est pas
surjective} ; la raison en est qu'il \uncertain{faut} prendre les \emph{duaux
topologiques} et non seulement
\note{« est injection » et « n'est pas surjective » sont soulignés, le
second doublement ; la phrase continue au feuillet suivant}

\page{15}
algébriques, \ill{} \uncertain{mettant} \uncertain{sur} $H^0(Y'; G) =
\struck{\ill{}} \prod_i \mathcal{O}_{x_i}$ la topologie \struck{produit
\uncertain{des} topologies} naturelle \ill{} l'\uncertain{anneau}
\uncertain{local}, [topologie \uncertain{fonctionnelle} \ill{}
\uncertain{produit}], et \uncertain{dans} $H^1(X'; F)$ \uncertain{une}
topologie \uncertain{quotient}, \uncertain{et} \uncertain{dans} \ill{}
\uncertain{discrète} \ill{} \ill{} $I$ \uncertain{pour} la topologie
\uncertain{quotient} \uncertain{sur} $\prod_i \mathcal{O}_{x_i}$ [N.B.\ $A$
\ill{} \ill{} \ill{} \ill{} \ill{} \ill{} \uncertain{complétion} \ill{} \ill{}
$\prod_i \mathcal{O}_{x_i}/A$ \ill{} \uncertain{séparé}].
\note{une incise cerclée, au-dessus de la cinquième ligne, ne se lit pas ;
tout le paragraphe est d'une écriture rapide}

\uncertain{Question} : \uncertain{dualité}\,…

\uncertain{Autres} \uncertain{questions} : $\operatorname{Ext}^i(K, A) = 0$
pour $i\,$ [$A$ anneau local \uncertain{noethérien} \uncertain{intègre}
complet, de \uncertain{corps} \uncertain{des fractions} $K$]
\note{le reste du feuillet est blanc}

\section{Notes antiques 1958 (Caractérisation des Cohen-Macaulay, des Gorenstein etc.)}
\note{titre de sa main, en haut à droite du feuillet 16, une chemise de
papier bistre par ailleurs vierge, qui ne reçoit pas de numéro de page ;
« Coh.\ Mac. » est abrégé sur la chemise}

\page{17}
\note{un plan, sur un feuillet coupé ; une grande accolade au crayon, dans la
marge gauche, embrasse les § 1 et 2}

\textbf{§ 1. \emph{Les foncteurs} $H^i_{\mathfrak{C}}$}

\begin{enumerate}
  \item Sur les foncteurs $A$-linéaires (de modules et modules
  \uncertain{topologiques}).
  \item Définition des $H^i_{\mathfrak{C}}(M)$, premières propriétés.
  \item Fonctorialité pour les $H^i_{\mathfrak{C}}(M)$ (y inclus
  cup-produit).
  \item Cas d'un anneau local $A$ [\struck{premières propriétés spéciales},
  \emph{caractérisation} des $H^i$ ; relations avec la codim.\ coh.].
  \item Rappels sur les Modules de Cohen-Macaulay.
  \item Résultats auxiliaires sur certains $\operatorname{Ext}^i_A(N, M)$
  [$M$ de C.-M., \ill{} supp.\ (\uncertain{par} $A$)].
  \item Cas particuliers.
\end{enumerate}
\note{« (y inclus cup-produit) » est cerclé ; le crochet de l'article 4 est
encadré, et une flèche le rattache à l'article 5, dont le numéro est
surchargé, comme celui de l'article 7}

\textbf{§ 2. \emph{Foncteurs et modules dualisants}} (théorie locale).
\emph{Le th.\ de dualité locale}.

\begin{enumerate}
  \item Définition et premières propriétés des foncteurs et modules
  dualisants.
  \item Extension du foncteur dualisant.
  \item Th.\ d'existence et unicité ; exemples.
  \item \struck{La théorie de dualité locale} Structure des $H^i(M)$.
  \item Module fondamental et \emph{th.\ de dualité locale}.
  \item Anneaux de Gorenstein.
\end{enumerate}

\textbf{§ 3. \emph{Les Foncteurs} $H_i$, \emph{et les complexes
dualisants}}

\page{18}
\struck{Remarque : \ill{} question de dualité, \ill{} \ill{} \ill{} la
validité de (iii) \ill{} \ill{} \ill{}, \ill{} \ill{} signifie que \ill{}
\ill{} $\dim A \leqslant \ill{}$\,\ill{}}
\struck{\ill{} donne \ill{}} \quad Théorème 4
\note{le début du feuillet est barré de longs traits obliques ; les mots
« Théorème 4 », encadrés, renvoient à l'énoncé plus bas}

\marginal{\uncertain{introduire} \ill{} \uncertain{définition} (iii)
\uncertain{et} certaines propriétés \ill{} de foncteurs $M \to
\operatorname{Hom}(M, \ill{})$}
\note{marginale écrite en oblique dans la marge gauche, en regard de la
définition, avec une flèche vers le « (ii) » ; lecture très incertaine}

\textbf{\emph{Définition.}} Soit $A$ un anneau \struck{local} noethérien
\uncertain{local} de dimension $n$. Un module $\Omega$ sur $A$ est dit un
« module fondamental » si il satisfait les deux conditions \struck{\ill{}
\ill{} \ill{}} suivantes :

(i) $\Omega$ est de type fini ;

(ii) \struck{\ill{}} $D(\Omega) \simeq H^n(A)$.

Un tel $\Omega$ est unique à \uncertain{isomorphisme} près \struck{si}
\uncertain{si} $\Omega = D(H^n(A))$ \struck{\ill{}} $=$ ($A$ étant local (ou
semi-local)), mais que \ill{}. On le note $\Omega_A$. D'autre part si $A$
est \uncertain{local} \emph{complet}, un tel $\Omega_A$ existe
(\uncertain{on} prend $D(H^n(A))$). Il existe \uncertain{aussi} et
\uncertain{équivalemment} si $A$ régulier : prendre $\Omega = A$\,!
\note{« local » est ajouté au-dessus de « complet » ; « et
équivalemment » est rattaché par une accolade à « si $A$ régulier »}

\struck{L'existence de $\Omega$ \ill{}} \struck{\ill{} \ill{}}
\struck{D'après le th.~3,} \struck{\ill{} \ill{} \ill{} \ill{} \ill{}}
\struck{Proposition. Soit $A$ un \ill{}\,…}
\struck{\ill{} $H^n(A) = \ill{}$}
\note{un bloc de lignes biffées et encadrées, entre la définition et le
théorème 4 ; une flèche le renvoie vers le haut}

Si $A$ \uncertain{est} \uncertain{cohen-macaulay}, alors module
\uncertain{fondamental} $=$ module dualisant.
\note{ajout encadré à droite, souligné}

\textbf{\emph{Théorème 4.}} Soit $A$ (\uncertain{local}) noethérien,
\uncertain{de dimension} $n$, muni d'un \struck{module dualisant}
\emph{\uncertain{fondamental}} $\Omega$. \uncertain{Le} \ill{} \ill{}
\uncertain{homomorphisme}
\[
  H^n(\Omega) \simeq H^n(A) \otimes \Omega \xrightarrow{\ \eta\ } I
\]
[\uncertain{déduit} de \uncertain{un} isomorphisme $H^n(A) \Rightarrow
D(\Omega)$], \uncertain{définit} \ill{} \uncertain{accouplement}
\uncertain{\ill{}}, \uncertain{qui} \uncertain{est} \ill{} \ill{}
$\Phi$\,\ill{} que \ill{} \uncertain{suite} \ill{}
\note{« local » est cerclé et ajouté au-dessus de « $A$ » ; le bas du
feuillet est d'une lecture de plus en plus incertaine}

\page{19}
\note{le feuillet entier est barré de deux longs traits obliques ; il se lit
néanmoins en grande partie}

\textbf{Complément au th.~3.} Si $A$ est de dimension $n$, [le foncteur $M
\mapsto H^n(M)$ \uncertain{est} exact à droite, et \struck{\ill{}}
\uncertain{il} \uncertain{commute} \ill{}
\[
  H^n(M) \xleftarrow{\ \sim\ } H^n(A) \otimes \underline{M}\,]
\]
\note{le crochet est de lui}

\textbf{Corollaire 2.} \struck{Corollaire 2 : si $D(\Omega) \simeq H^n(A)$,
\ill{} \ill{} \ill{}} \struck{\ill{} \ill{}}
\uncertain{et} \uncertain{si} \uncertain{on} a $\operatorname{Ext}^i(k, M)
\neq 0$, \uncertain{alors} \ill{} $\operatorname{Ext}^j(k, M) = 0$ pour $j >
i$ \ill{}, $\operatorname{Ext}^q(k, M) \neq 0$ …
\note{la ligne du corollaire, biffée d'un trait, porte au-dessus des
reprises qui ne s'accordent pas ; à droite, biffé : « $H^n(\Omega) \simeq
D(\operatorname{Hom}(\Omega, \Omega))$ »}

Soit \uncertain{donc} $\Omega$ \uncertain{dualisant}, \uncertain{et}
\[
  H^n(\Omega) \simeq \struck{\ill{}}\, H^n(A) \otimes \Omega
  \xrightarrow{\ \varphi\ } I
\]
\struck{\ill{} \ill{}} Je \uncertain{suppose} \ill{} \ill{} \ill{} \ill{}
$\Omega \xrightarrow{\varphi_2} D(H^n(A))$
\[
  \ill{}\quad H^n(A) \xrightarrow{\varphi_1} D(\Omega).
  \quad \uncertain{Un} \ill{} \ill{}
\]
\uncertain{Accouplements}
\[
  H^p(M) \times \operatorname{Ext}^{n-p}(M, \Omega) \longrightarrow
  H^n(\Omega) \rightsquigarrow I
\]
d'où
\[
  H^p(M) \longrightarrow D(\operatorname{Ext}^{n-p}(M, \Omega))
\]
\uncertain{Explicitons} \ill{} \ill{} \uncertain{dimension} cohomologique
\[
  H^n(M) \simeq \struck{\ill{}}\, H^n(A) \otimes M \longrightarrow
  D(\operatorname{Hom}(M, \Omega)) \simeq D(\Omega) \otimes M
\]
\struck{\ill{} \ill{}} \uncertain{précisément} : $\varphi_1 \otimes
\operatorname{id}_M$ ; \uncertain{pour} \uncertain{que} \ill{}
\uncertain{isomorphisme} pour \uncertain{tout} $M$, il \uncertain{faut}
\uncertain{et} \uncertain{il suffit} \ill{} \uncertain{soit} \ill{}
\uncertain{isomorphisme}. \uncertain{Considérons} le \ill{} $M = \Omega$,
\ill{} \uncertain{donne}
\[
  \struck{H^n(\Omega) \ill{}(A)}
\]
\[
  H^n(\Omega) \times \widehat{\operatorname{Hom}(\Omega, \Omega)}
  \longrightarrow H^n(\Omega) \longrightarrow I
\]
d'où
\[
  \operatorname{Hom}(\Omega, \Omega) \xrightarrow{\ \sim\ } D(H^n(\Omega))
\]
\uncertain{transformant} l'identité en $\varphi$. \uncertain{On}
\uncertain{dira} \struck{\ill{} \ill{} \ill{} \ill{} \ill{} :}

\struck{(i) $H^n(\Omega)$ \ill{} dualisant \ill{}}

\struck{(ii) $\operatorname{Hom}(\Omega, \Omega) \simeq A$} \quad (iii)
\marginal{\uncertain{utiles} \uncertain{tous} ?}
\note{la marginale, dans la marge gauche, porte un trait vertical qui
embrasse le bas du feuillet ; le « (iii) » final est relié par un trait
sinueux aux lignes biffées (i) et (ii)}

\page{20}
\uncertain{isomorphisme} pour les \uncertain{dimensions} $p \geqslant k$
($k$ \uncertain{entier} \uncertain{donné} $\leqslant n$), il faut et il
suffit \uncertain{qu'on} \uncertain{ait}
\[
  H^i(A) = 0 \quad \text{pour } k \leqslant i < n.
\]
\note{ce premier alinéa, qui continue le feuillet précédent, est marqué
d'un trait vertical dans la marge}

\textbf{\emph{Corollaire 1.}} \uncertain{Pour} que \ill{} \ill{} \uncertain{soit}
\uncertain{un} \uncertain{isomorphisme} pour \struck{\ill{}} \uncertain{tout}
$p$, il \uncertain{faut} et il suffit que \uncertain{les}
\uncertain{composantes} \uncertain{de} $A$ \uncertain{soient}
\struck{\ill{}} \uncertain{anneaux} de Coh.\ Mac.\ de dimension $n$.
\note{« Composantes » est ajouté au-dessus de la ligne ; un mot biffé et
souligné d'une ondulation précède « anneaux »}

\struck{\ill{}} \textbf{\emph{Corollaire 2}} \struck{Proposition~2}. Soit $A$
\uncertain{un} anneau \uncertain{local} de Coh.\ Mac., $\Omega$ un module
\struck{dualisant} \uncertain{fondamental} \ill{} \uncertain{sur} $A$.
\uncertain{Alors} \ill{} :
\[
\left\{
\begin{aligned}
  &\operatorname{Ext}^i(k, \Omega) = 0 \quad \text{si } i \neq n\\
  &\operatorname{Ext}^n(k, \Omega) \simeq k
\end{aligned}
\right.
\]
\note{« Corollaire 2 » est écrit au-dessus de « Proposition 2 », biffé ;
« loc. » est ajouté au-dessus de la ligne, « fond. » au-dessus de
« dualisant », biffé}

\uncertain{En} \uncertain{effet}, $\operatorname{Ext}^i(k, \Omega)$ est
\uncertain{dual} de $H^{n-i}(k)$, \uncertain{qui} est \uncertain{nul} si
$n - i \neq 0$, et \uncertain{qui} est $\simeq \underline{k}$ si $n - i =
0$.
\note{une accolade dans la marge embrasse ces deux lignes}

\uncertain{En} \uncertain{particulier}

\textbf{\emph{Corollaire 3.}} $A$, $\Omega$ \uncertain{comme} \ill{}
\uncertain{ci-dessus}. \uncertain{Alors} $\Omega$ est un module de
Coh.\ Mac.\ (\uncertain{sur} $A$, \uncertain{de} \uncertain{dimension}
$n$), \uncertain{et} \ill{} :
\[
\left\{
\begin{aligned}
  &H^i(\Omega_{\ill{}}) = 0 \quad \text{si } i \neq n\\
  &H^n(\Omega_{\ill{}}) \text{ \uncertain{est} \uncertain{dualisant}}
\end{aligned}
\right.
\]
De plus, \struck{$\Omega_A$ \uncertain{est} \ill{} \ill{} fidèle}
\[
\left\{
\begin{aligned}
  &\operatorname{Ext}^i(\Omega, \Omega) = 0 \quad \text{si } i \neq 0\\
  &\struck{\operatorname{Ext}}\ \operatorname{Hom}(\Omega, \Omega) \simeq A
\end{aligned}
\right.
\]
\note{l'indice de $\Omega$ dans les deux lignes de la première accolade est
surchargé : peut-être un astérisque, peut-être un « $A$ »}

\end{document}
